\documentclass[11pt, a4paper]{common/document_template}
\usepackage{common}

\begin{document}

\begin{titlepage}
  \begin{center}
	\Huge{Летние конференции Турнира городов}
  \end{center}
\end{titlepage}
 
\chapter{2019. Номер 31}

\begin{exercise}
В некоторых клетках прямоугольной таблицы нарисованы звездочки. Известно, что для любой отмеченной клетки количество звездочек в её столбце совпадает с количеством звездочек в её строке. Докажите, что число строк в таблице, в которых есть хоть одна звездочка, равно числу столбцов таблицы, в которых есть хоть одна звездочка.
\end{exercise}

\begin{exercise}
В прямоугольной таблице $m$ строк и $n$ столбцов, где $m < n$. В некоторых клетках таблицы стоят звёздочки, так что в каждом столбце стоит хотя бы одна звёздочка. Докажите, что существует хотя бы одна такая звёздочка, что в одной строке с нею находится больше звёздочек, чем с нею в одном столбце.
\end{exercise}

\begin{exercise}
В библиотеке на полках стоят книги, ровно $k$ полок пусты. Книги переставили так, что теперь пустых полок нет. Докажите, что найдётся хотя бы $k + 1$ книга, которая теперь стоит на полке с меньшим числом книг, чем стояла раньше.
\end{exercise}

\begin{exercise}
Дан выпуклый многогранник, у которого нет четырехугольных и пятиугольных граней. Докажите, что у него по крайней мере 4 треугольные грани.
\end{exercise}

\begin{exercise}
Граф, изображенный на плоскости, называется квазипланарным, если у него не найдется трёх ребер, попарно пересекающихся во внутренних точках. Докажите, что число рёбер в квазипланарном графе на $n$ вершинах не превосходит $8n - 20$.
\end{exercise}
 
\begin{exercise}
Граф называется спичечным, если он может быть изображён на плоскости так, чтобы все его ребра являлись непересекающимися друг с другом отрезками длины 1. Докажите, что не существует спичечного графа, степень каждой вершины которого равна 5.
\end{exercise} 
 
\end{document}